\begin{satz}{Injektivität linearer Operatoren und Operatornorm der Inversen}
            {InjektivitaetLinearerOperatorenUndOperatornormDerInversen}
Seien $(E,\norm\cdot_E),(F,\norm\cdot_F)$ $\K$--Banachräume, $T\in L(E,F)$
injektiv und $T(E)$ sei abgeschlossen in $F$.\\
Dann gilt \[\inf\menge{\norm{Tx}_F}{x \in\kran E 0 1}>0\text.\]
Weiter ist $T\inv$ eine Abbildung, $T\inv \in L(T(E), E)$
und es gilt \[
\opnorm{T\inv} {T(E)} E 
   = \frac 1 {\inf\menge{\norm{Tx}_F}{x \in\kran E 0 1}}\text.
\]
\end{satz}
\begin{proof}
Sei $G\coloneqq T(E)$.\\
Dann ist $G$ wegen der Linearität von $T$ ein Vektorraum.\\
Offenbar ist $\norm\cdot_G\coloneqq \restrict{\norm\cdot_F}G$ eine Norm auf $G$.\\
Da $G$ nach Voraussetzung abgeschlossen in $F$ ist, ist $G$ als abgeschlossene
Menge eines vollständigen metrischen Raums selbst vollständig und damit
insgesamt ein Banachraum.\\
Sei $S : E \rightarrow G, x \mapsto Tx$.\\
Dann ist $S$ bijektiv (weil $T$ injektiv ist), insbesondere surjektiv.\\
Offenbar ist auch $S$ in dieser Normwahl stetig.\\
Nach dem Satz von der offenen Abbildung ist damit $S$ offen.\\
Nun ist $S$ aber auch invertierbar, sodass damit auch $S\inv$ stetig ist.\\
Da $T$ als Relation links-- und rechtseindeutig ist, ist es auch $T\inv$.\\
Damit ist $T\inv$ eine injektive Abbildung von $T(E)$ nach $E$ und stimmt als
solche mit $S\inv$ überein, ist also inbesondere linear und stetig.\\
Seien $c\coloneqq  \opnorm {S\inv} G E$ und $x \in \kran E 0 1$.\\
Dann ist $c > 0$ und es gilt:
\[1 =\norm x _E = \norm{S\inv(Sx)}_E \le \opnorm{S\inv} G E \norm{Tx}_G 
             = c \norm{Tx}_G\text.\]
Also ist $\frac 1 c$ eine untere Schranke der Menge
$\menge{\norm{Tz}_F}{z\in\kran E01}$, sodass auch
$\inf\menge{\norm{Tz}_F}{z\in\kran E01} \ge \frac 1 c > 0$ gilt.\\
Offenbar ist 
$\menge{\frac{Tz}{\norm{Tz}_F}}{z \in \kran E 0 1} \subseteq \kran G01$,
während für $w \in \kran G01$ mit der Setzung 
$v\coloneqq \frac{S\inv w}{\norm{S\inv w}_E}$ gilt:
$\norm v_E = 1$ und
\[\frac{Tv}{\norm{Tv}_F} = \frac{T\left(\frac{S\inv w}{\norm{S\inv w}_E} \right)}{\norma{T\left(\frac{S\inv w}{\norm{S\inv w}_E}\right)}_F}
= \frac{\frac 1 {\norm{S\inv w}_E} w}{\frac{\norm w_F}{\norm{S\inv w}_E}}
\overset{\norm w _F=1} = w\text.\]
Somit ist
$\mengea{\frac{Tz}{\norm{Tz}_F}}{z \in \kran E 0 1} = \kran G01$, also:
\begin{align*}
\opnorm {T\inv}{T(E)}E & = \opnorm {S\inv} G E 
        = \sup\menge{\norm{S\inv z}_E}{z \in \kran G 0 1}\\
        &= \sup\mengea{\norma{S\inv \left(
                  \frac{Ty}{\norm{Ty}_F}
                  \right)}_E}{y \in \kran E 0 1}\\
        &= \sup\mengea{\frac 1 {\norm {Ty}_F}}{y \in \kran E 0 1}\\
        &= \frac 1 {\inf\menge{\norm {Ty}_F}{y \in \kran E 0 1}}\text.
\end{align*}
\end{proof}